\documentclass[11pt]{article}
\usepackage[margin=1in]{geometry}
\usepackage[T1]{fontenc}
\usepackage[utf8]{inputenc}
\usepackage{lmodern}
\usepackage[expansion=false]{microtype}

\pagestyle{empty}

\begin{document}

\noindent Dear Editors,

\vspace{0.75\baselineskip}

\noindent Please consider my manuscript, ``Protective Computing: An Architectural Framework for Safety-Critical Systems Under Conditions of Human Vulnerability,'' for publication as a Research Article.

\vspace{0.75\baselineskip}

\noindent This paper introduces Protective Computing as a systems-engineering orientation designed for contexts in which human--computer interaction occurs under vulnerability (medical crisis, coercion, infrastructural instability, or cognitive degradation). The manuscript formalizes Stability Bias ($H = f(D, V, I)$), defines a Vulnerability State Machine ($S_0$--$S_3$), specifies five normative design principles using RFC-style language, and proposes a provisional evaluation metric (Protective Legitimacy Score).

\vspace{0.75\baselineskip}

\noindent The paper synthesizes and extends work in safety-critical HCI, resilience engineering, local-first architecture, privacy-by-design, and digital safety research. It advances the claim that architectural legitimacy in safety-critical domains must be evaluated by system behavior under degraded human and infrastructural conditions rather than ideal stability assumptions.

\vspace{0.75\baselineskip}

\noindent A reference implementation (PainTracker) is provided to demonstrate technical feasibility. The implementation evidence is presented conservatively and is not claimed as large-scale validation.

\vspace{0.75\baselineskip}

\noindent The manuscript has not been published elsewhere and is not under consideration by another journal. There are no conflicts of interest to declare.

\vspace{0.75\baselineskip}

\noindent I believe this work will be of interest to researchers in HCI, trustworthy systems, digital medicine, and safety-critical computing.

\vspace{0.75\baselineskip}

\noindent Thank you for your consideration.

\vspace{1.25\baselineskip}

\noindent Sincerely,\\
K.\ Overton

\end{document}
